\notechapter{N-20220810}{Basic Inequalities, Cauchy, Triangle Inequality, and Coordinate Geometry}

\section{Basic inequalities}

Every square is nonnegative.  For real numbers,
\[
  (a-b)^2\ge0,
\]
so
\[
  a^2+b^2\ge2ab,
\]
with equality iff $a=b$.  For three variables,
\[
  a^2+b^2+c^2\ge ab+bc+ca,
\]
because
\[
  2(a^2+b^2+c^2-ab-bc-ca)=(a-b)^2+(b-c)^2+(c-a)^2\ge0.
\]
Equality occurs exactly when all variables involved in the squared differences coincide.

\section{AM--GM}

For $a,b\ge0$,
\[
  \frac{a+b}{2}\ge\sqrt{ab},
\]
with equality iff $a=b$.  Equivalently,
\[
  a+b\ge2\sqrt{ab}.
\]
The note also records the three-variable form
\[
  \frac{a+b+c}{3}\ge\sqrt[3]{abc},\qquad a,b,c\ge0,
\]
and uses it repeatedly in examples.

A typical application is to expressions such as
\[
  \frac{a+b+c}{3}\ge\sqrt[3]{abc}
  \quad\Longrightarrow\quad
  a+b+c\ge3\sqrt[3]{abc}.
\]
The useful habit is to look for terms whose product is fixed.  AM--GM then turns a fixed product into a lower bound for a sum.

\section{Cauchy's inequality}

Cauchy's inequality has two useful forms.  The vector form is
\[
  (a_1b_1+\cdots+a_nb_n)^2
  \le
  (a_1^2+\cdots+a_n^2)(b_1^2+\cdots+b_n^2).
\]
The Engel or Titu Andreescu form is
\[
  \frac{x_1^2}{a_1}+\cdots+\frac{x_n^2}{a_n}
  \ge
  \frac{(x_1+\cdots+x_n)^2}{a_1+\cdots+a_n},\qquad a_i>0.
\]
This second form follows by applying Cauchy to
\[
  \left(\sum_i \frac{x_i^2}{a_i}\right)
  \left(\sum_i a_i\right)
  \ge
  \left(\sum_i x_i\right)^2.
\]

\section{Triangle inequality and its variants}

The page then discusses absolute values:
\[
  |a+b|\le |a|+|b|,\qquad
  |a-b|\le |a-c|+|c-b|.
\]
The second formula is the triangle inequality with an intermediate point $c$.  The note also uses the reverse triangle inequality
\[
  ||a|-|b||\le |a-b|.
\]
A useful proof is
\[
  |a|=|a-b+b|\le |a-b|+|b|,\qquad
  |b|=|b-a+a|\le |a-b|+|a|.
\]
Together these imply the reverse inequality.

Absolute-value inequalities are not merely algebraic sign manipulations; they encode distances on the real line.

\section{Worked AM--GM examples}

One example asks for a lower bound of expressions of the form
\[
  x+\frac1x,\qquad x>0.
\]
By AM--GM,
\[
  x+\frac1x\ge2.
\]
Similarly,
\[
  x+y+z\ge3\sqrt[3]{xyz}
\]
when $x,y,z>0$.  If a problem fixes $xyz$, this gives a direct lower bound for the sum.

Another common pattern is
\[
  \frac{x}{y}+\frac{y}{z}+\frac{z}{x}\ge3,\qquad x,y,z>0,
\]
because the product of the three terms is $1$.

\section{A Cauchy example}

For positive $a,b,c$, a typical inequality is
\[
  \frac{a^2}{b+c}+\frac{b^2}{c+a}+\frac{c^2}{a+b}
  \ge
  \frac{(a+b+c)^2}{2(a+b+c)}
  =\frac{a+b+c}{2}.
\]
This is exactly Engel form of Cauchy.  It turns a sum of fractions into one square divided by one denominator.

\section{Coordinate geometry at the end of the note}

The last page shifts to coordinate geometry.  The diagrams show points, vectors, and lines, with formulas for distances and slopes.  The essential formulas are:
\[
  AB=\sqrt{(x_A-x_B)^2+(y_A-y_B)^2},
\]
\[
  m=\frac{y_2-y_1}{x_2-x_1},
\]
and the point-slope form of a line
\[
  y-y_0=m(x-x_0).
\]

A recurring idea is that geometric claims can be turned into algebraic equalities.  For example, perpendicularity of two nonvertical lines is
\[
  m_1m_2=-1,
\]
while parallelism is
\[
  m_1=m_2.
\]
Distance and slope are the coordinate versions of length and direction.


\section{What the examples reveal: three structures}

Choosing the right inequality language is the main skill.  Squares prove elementary inequalities, AM--GM handles fixed products, Cauchy handles dot products and sums of fractions, and triangle inequalities handle distances.  The advanced form of the same habit is to identify the structure before calculating.

\begin{center}
\small
\begin{tabularx}{\textwidth}{@{}lYY@{}}
\toprule
Elementary inequality & Higher structure & Equality meaning\\
\midrule
Cauchy--Schwarz & inner product / projection / Gram matrix & vectors are dependent\\
Triangle inequality & normed vector space & same direction in strictly convex cases\\
AM--GM & convexity or concavity of \(\log\) & all variables equal\\
Positive square & positive semidefinite cone & quadratic form degenerates\\
\bottomrule
\end{tabularx}
\end{center}

\section{Cauchy--Schwarz and Gram determinants}

Cauchy--Schwarz says
\[
  |\langle x,y\rangle|\le \|x\|\|y\|.
\]
It follows from the nonnegativity of
\[
  \|x-ty\|^2\ge0
\]
or, equivalently, from positivity of the Gram matrix
\[
  G(x,y)=
  \begin{pmatrix}
  \langle x,x\rangle&\langle x,y\rangle\\
  \langle y,x\rangle&\langle y,y\rangle
  \end{pmatrix}.
\]
The inequality \(\det G(x,y)\ge0\) is exactly Cauchy--Schwarz.  Equality means \(G\) has rank one, i.e. the two vectors are linearly dependent.

In Hilbert spaces, the same inequality controls Fourier coefficients:
\[
  |\langle f,e_n\rangle|\le \|f\|_2\|e_n\|_2.
\]
This is why an inequality first met in finite sums becomes indispensable in functional analysis.

\section{Triangle inequality and normed spaces}

The triangle inequality
\[
  \|x+y\|\le \|x\|+\|y\|
\]
defines the geometry of a normed vector space.  The usual absolute value, Euclidean length, sup norm, and \(L^p\) norms are all instances.  Minkowski's inequality is precisely the triangle inequality in \(L^p\):
\[
  \|f+g\|_p\le \|f\|_p+\|g\|_p,
  \qquad
  \|f\|_p=\left(\int |f|^p\right)^{1/p}.
\]
Thus a contest-level inequality about square roots and sums is structurally the same statement as the triangle inequality for functions.

\section{AM--GM, convexity, and cones}

For positive variables, AM--GM follows from concavity of \(\log\):
\[
  \log\left(\frac1n\sum_i x_i\right)\ge \frac1n\sum_i\log x_i.
\]
The equality condition says all variables are equal, meaning there is no spread.

Many inequalities also describe membership in a convex cone.  For instance,
\[
  A\succeq0\quad\Longleftrightarrow\quad x^TAx\ge0\text{ for all }x
\]
defines the cone of positive semidefinite matrices.  Basic ``square is nonnegative'' arguments are the one-dimensional shadow of this cone viewpoint.
